\documentclass[a4paper,12pt]{article}

\usepackage[utf8]{inputenc} % encodé en utf-8
\usepackage[T1]{fontenc} % compatible avec les accents

\usepackage[french]{babel} % rédigé en français
\usepackage[hyphens]{url} % formatte les liens en autorisant la césure au niveau des traits d'union
\usepackage[urlcolor=black,colorlinks=true,linkcolor=black,citecolor=black]{hyperref} % liens cliquables mais non colorés

\usepackage[top=25mm, bottom=25mm, left=35mm, right=25mm]{geometry} % définit les marges
\usepackage{array} % permet l'alignement vertical dans les tableaux
\usepackage{setspace} % permet de régler l'interligne
\onehalfspacing % définit l'interligne 1,5 par défaut

\usepackage{graphicx} % gestion des images
\graphicspath{ {../images/}}
\usepackage{array} % gestion des tableaux
\usepackage{csquotes} % gestion des guillemets

\usepackage{fourier} % utilise une autre police que celle par défaut (Computer Modern) fourier - venturis - libertine - tgtermes - kpfonts - lmodern - txfonts - newtxtext,newtxmath - mathptmx - stix 

\usepackage[backend = biber, style = authoryear, useprefix = true, maxbibnames = 99]{biblatex}


%\addbibresource{}

\DefineBibliographyExtras{french}{%
\protected\def\mkbibnamefamily#1{%
\textsc{\textnohyphenation{#1}}}
}%

\DefineBibliographyStrings{french}{%
    page  = {\lowercase{p.}}, % for single page number
    pages = {\lowercase{pp.}}, % for multiple page numbers
}

\DefineBibliographyExtras{french}{\renewcommand*{\bibrangedash}{\textendash}} % Pour rendre les "--" correctement

% \renewcommand*{\mkibid}{\emph} %ibid en italique  if authoryear-ibid

\usepackage{float}

\setcounter{secnumdepth}{3} %Ne compte pas plus loin que 1.1.1
\setcounter{tocdepth}{3}

\usepackage{enumitem} %% utilisé pour la présentation de fausse table des matières dans la table des matières commentée

\newlist{legal}{enumerate}{10}
\setlist[legal]{label*=\arabic*.}

\usepackage{ragged2e} % Permet de placer plus facilement les tableaux

\newcolumntype{R}[1]{>{\RaggedRight}p{#1}}

\title{Projet STIC-B505 : Budgetsquirrel, une application de gestion de budget} 
\author{Milena Albu \and Marie Huysman} % vos prénom et nom
\date{} % pas de date

%\usepackage{fancyhdr}
%\pagestyle{fancy} % Turn on the style
%\fancyhf{}
%\fancyhead[L]{\footnotesize \nouppercase \leftmark}
%\fancyhead[R]{\footnotesize \nouppercase \rightmark}
%\fancyfoot[C]{\thepage}
%
%\renewcommand{\chaptermark}[1]{%
%\markboth{#1}{}}
%
%    \renewcommand{\sectionmark}[1]{\markright{\thesection \ #1}{}}

\usepackage{longtable}


\begin{document} % début du corps du texte
\sloppy

\begin{titlepage} % style spécial pour la page de garde
\singlespacing % remet l'interligne normal pour cette page

%%alternative à la bannière image
%\begin{center} 
%\huge{\textsc{Université Libre de Bruxelles}}\\
%\vspace{.5cm}
%\LARGE{Faculté de Lettres, Traduction et Communication}
%\end{center}
%% Bannière image
\begin{figure}[h]
\noindent
\makebox[\textwidth]{\includegraphics[scale=0.6]{banner.png}}
\end{figure}

\vfill % laisse LaTeX gérer l'espacement vertical

%VOIR SI LE TITRE CONVIENT

\begin{center}
\LARGE{Conception et gestion de banques de données : Rapport de projet}\\ % ajoutez votre titre ici
\Large{ Budget Squirrel : Application de gestion de budget}
\vspace{.5cm}
\end{center}

\vfill

\begin{tabular}{b{5.5cm}b{7.5cm}} % taille et alignement des cellules du tableau
 \\ Milena \textsc{Albu} \\ Marie \textsc{Huysman} \\ & Rapport de projet écrit dans le cadre du cours STIC-B505 : Conception et gestion de banques de données \\
\end{tabular}

\vfill

\begin{center}
Année académique 2020--2021 % vérifiez que l'année est correcte !
\end{center}

\end{titlepage}

\tableofcontents

\thispagestyle{plain}

\newpage

\section{Introduction}

Si besoin de citer : \parencite[2--4]{SOURCE}
Notes de bas de page \footnote{informations suppl.}

Dans le cadre de ce projet, nous avons choisi de développer une application permettant la gestion de budget à une échelle personnelle.
Après une brève présentation du domaine d'application, nous présenterons le schéma conceptuel, ainsi que sa traduction en schéma relationnel.
Ensuite, nous décrirons les différents aspects des scripts SQL que nous avons mis en place [...] .
Enfin, nous exposerons brièvement l'application Web que nous avons développée. 
Vous trouverez, en annexe à ce rapport, deux scripts sql : \verb\initdb.sql\, qui contient tous les scripts permettant la création de  la base de donnée en elle-même, et \verb\data.sql\, qui peuple la base de données de quelques informations, permettant ainsi de parcourir les différentes pages de l'application.

\subsection{Domaine d'application}

Utilisateurs cibles, cas d'utilisation du système

\newpage

\section{Schéma conceptuel}

Schéma conceptuel documenté, accompagné d'une description et justification des choix

\newpage 

\section{Schéma logique relationnel}

Schéma logique relationnel traduisant le schéma conceptuel, avec une explication du choix des clés de traductions (héritage, associations, contraintes...

\newpage

\section{Code SQL}

\subsection{Description du code de création de la base de données}

Description du code sql de création de la db 
!! La création de la db doit : créer des tables , avoir des contraintes de clé primaire, unique et externe, avoir des contraintes sous forme de prédicat de table et de colonne, avoir des procédures stockées et des déclencheurs

\subsection{Éléments de SQL avancé}

\subsubsection{Contraintes garanties par des check}

\subsubsection{Requêtes de consultation}

\subsubsection{Requêtes de mise à jour}


\subsubsection{Vues utilisées}

Ce qu'on utilise pour l'historique et les différentes statistiques

\subsubsection{Déclencheurs}

Triggers à utiliser, sans doute, pour le calcul des bilans, et autres 

\subsubsection{Droits d'accès}

Voir si les droits d'accès peuvent être utilisés comme dans le rapport publié par Servais, pour l'instant ça reste un accès root défini par une session > problématique?



\newpage
\section{Description de l'application Web}

Description brève des pages de l'application web et de leur fonctionnement php

\newpage

\section{Développement du projet}

Quelques mots sur la logique du travail : paired programming, outils employés, défis rencontrés, ...

\newpage
\section{Conclusion}

Les points forts et faiblesses d’application, le temps nécessaire pour développer le concept avec deux personnes, les sources consultées pour la partie pratique (i.e. les slides, ou techniques vues au TP, ainsi que les conseils d’assistante, mais aussi w3schools ou bien stackoverflow).
Finalement quelques mots sur l'expérience du projet.




\newpage

%voir si bibliographie nécessaire
%\chapter*{Bibliographie} 
%\addcontentsline{toc}{chapter}{Bibliographie}
%\markboth{Bibliographie}{}
%
%\printbibliography[heading=none]
%
%\fussy

%Voir si annexes nécessaire
%\chapter*{Annexes}
%\addcontentsline{toc}{chapter}{Annexes}
%\markboth{Annexes}{}

\end{document} % fin du corps du texte